\documentclass[11pt]{article}
\usepackage[T1]{fontenc}
\usepackage{lmodern}
\usepackage{microtype}
\usepackage{amsmath,amssymb,amsthm,mathtools}
\usepackage[margin=1in]{geometry}
\usepackage{xcolor}
\usepackage{booktabs,array}
\usepackage[colorlinks=true,linkcolor=blue!55!black,citecolor=blue!55!black,urlcolor=blue!55!black]{hyperref}
\usepackage{enumitem}

\newtheorem{theorem}{Theorem}
\newtheorem{lemma}{Lemma}
\newcommand{\C}{\mathbb C}
\newcommand{\ran}{\operatorname{ran}}
\newcommand{\rank}{\operatorname{rank}}
\newcommand{\id}{\operatorname{id}}
\newcommand{\M}{M_3(\C)}

\definecolor{statusbg}{RGB}{239,246,255}
\definecolor{statusrule}{RGB}{37,99,235}

\title{The q-pure-map conjecture on $M_3(\C)$\\
\large A literature-implied full solution of the $n=3$ subcase}
\author{Automated Mathematics Run \texttt{fa\_banach\_001}, lane 2}
\date{27 August 2026}

\begin{document}
\maketitle

\noindent\colorbox{statusbg}{\parbox{0.94\textwidth}{%
\textbf{Status.} Literature-implied answer (full $n=3$ subcase), likely valid.
The implication assembled below was not found stated explicitly in the cited papers.
The conjecture for $n\geq4$ remains outside the scope of this note.}}

\section{Source conjecture and exact scope}

In the concluding paragraph of arXiv:0905.2708v2, PDF page 37,
Jankowski asks for all unital q-pure maps $\phi:M_n(\C)\to M_n(\C)$ and
states the conjecture that every such map is invertible or has rank one
\cite{Jankowski2009}.  The same paragraph notes that the $n=2$ case is
proved in later work and suggests using the limit map for $n=3$.

This note proves exactly that next case.

\begin{theorem}[The $n=3$ case]\label{thm:main}
If $\phi:M_3(\C)\to M_3(\C)$ is unital and q-pure, then either
$\rank\phi=1$ or $\phi$ is invertible.
\end{theorem}

Here q-pure has the source-paper meaning: the q-subordinates of $\phi$ are
precisely $0$ and
\[
  \phi^{(s)}=\phi(I+s\phi)^{-1},\qquad s\geq0.
\]

\section{Identification mechanism}

The proof joins two results that were developed for different purposes.
The 2010 paper classifies every possible limit idempotent
\[
  L_\phi=\lim_{s\to\infty}s\phi(I+s\phi)^{-1}
\]
for a unital q-positive map on $M_3(\C)$, and proves that a q-positive map
annihilating a nonzero positive matrix is not q-pure
\cite[Theorem 3.9, Corollary 3.10, Proposition 4.5]{Jankowski2010}.
The 2018 paper proves that the range of a finite-dimensional, index-zero,
q-pure q-weight is a factor under its Choi--Effros product
\cite[Theorem 5.7]{JMP2018}.

The bridge is supplied by the boundary-weight double already constructed in
the source paper.  The next lemma records the bridge explicitly because this
is the nontrivial identification on which the result depends.

\begin{lemma}[Boundary-weight bridge]\label{lem:bridge}
Let $\phi:M_n(\C)\to M_n(\C)$ be unital and q-pure.  Choose a normalized
unbounded vector boundary weight $\nu$ of the form required in Lemma 4.3 of
\cite{Jankowski2009}, and let $\omega$ be the q-weight dual to the boundary
weight map of the double $(\phi,\nu)$. Then:
\begin{enumerate}[label=(\roman*),leftmargin=2.2em]
  \item $\omega$ is q-pure and has index zero;
  \item $\ran\omega=\ran\phi$;
  \item the Choi--Effros limit of $\omega$ is $L_\phi$.
\end{enumerate}
Consequently $\ran\phi$, with product $A\star B=L_\phi(AB)$, is a factor.
\end{lemma}

\begin{proof}
The source paper defines the dual q-weight by
\[
  \omega(A)=\phi\bigl(\Omega_\nu(A)\bigr),
\]
where $\Omega_\nu$ applies $\nu$ entrywise.  Lemma 4.3 of
\cite{Jankowski2009} says that every subordinate CP-flow of the flow induced
by $(\phi,\nu)$ is induced by $(\psi,\nu)$ for a q-subordinate
$\psi\leq_q\phi$.  Since q-purity of $\phi$ makes all such $\psi$ belong to
the chain $\{0\}\cup\{\phi^{(s)}:s\geq0\}$, the q-subordinates of $\omega$
are totally ordered.  Thus $\omega$ is q-pure in the sense of
\cite[Definition 1.2]{JMP2018}.  Corollary 3.3 of
\cite{Jankowski2009} gives index zero.

Normalization gives $\nu(I-\Lambda(1))=1$.  For every $B=(b_{ij})\in
M_n(\C)$, the matrix $A=(b_{ij}(I-\Lambda(1)))$ belongs to the null boundary
algebra and satisfies $\Omega_\nu(A)=B$.  Hence $\Omega_\nu$ is onto and
$\ran\omega=\ran\phi$.

For $t>0$, put $s_t=\nu_t(\Lambda(1))$.  Equation (3) in
\cite{Jankowski2009} gives the generalized boundary representation
\[
  \pi_t^\#=\phi(I+s_t\phi)^{-1}\Omega_{\nu_t}.
\]
Since $\Omega_{\nu_t}(\Lambda(B))=s_tB$,
\[
  \pi_t^\#\Lambda(B)=s_t\phi(I+s_t\phi)^{-1}(B).
\]
Unboundedness of $\nu$ implies $s_t\to\infty$ as $t\downarrow0$, so the
right-hand side converges to $L_\phi(B)$.  This proves (iii).
The final assertion is now exactly Theorem 5.7 of \cite{JMP2018}.
\end{proof}

\section{Proof of Theorem \ref{thm:main}}

\begin{proof}
For a unital q-positive $\phi$ on $M_3(\C)$, the 2010 limit-map lemma shows
that $L_\phi$ is a unital completely positive idempotent and that
\[
  \ran L_\phi=\ran\phi,
  \qquad
  \ker L_\phi=\ker\phi.
\]
Theorem 3.9 of \cite{Jankowski2010} classifies such idempotents up to unitary
conjugacy.  Proposition 4.5 of the same paper excludes every case in which
$L_\phi$ annihilates a nonzero positive matrix: equality of kernels would
make $\phi$ annihilate that matrix, contradicting q-purity.

Thus, apart from a faithful rank-one map, the only possible classified forms
for $L_\phi$ are the following four forms (the labels are those of
\cite{Jankowski2010}):
\[
\begin{array}{ccl}
\text{(IV)}&L(A)=\operatorname{diag}(a_{11},a_{22},a_{33}),
  &\ran L\cong\C^3,\\[2pt]
\text{(V)}&L(A)=a_{11}\oplus
  \begin{pmatrix}a_{22}&a_{23}\\a_{32}&a_{33}\end{pmatrix},
  &\ran L\cong\C\oplus M_2(\C),\\[7pt]
\text{(VI)}&L(A)=a_{11}\oplus bI_2,
  \quad b=\lambda a_{22}+(1-\lambda)a_{33},
  &\ran L\cong\C\oplus\C,\\[2pt]
\text{(VII)}&L(A)=A,
  &\ran L=M_3(\C).
\end{array}
\]
In cases (IV)--(VI), the range is already closed under ordinary matrix
multiplication and $L$ fixes every range element.  Therefore its
Choi--Effros product $X\star Y=L(XY)$ is ordinary multiplication.  Each of
$\C^3$, $\C\oplus M_2(\C)$, and $\C\oplus\C$ has nontrivial center, so none
is a factor.  Lemma \ref{lem:bridge} eliminates these three cases.

If $L_\phi$ is the faithful rank-one form, equality of ranges gives
$\rank\phi=1$.  The only remaining case is (VII), where
$\ran\phi=\ran L_\phi=M_3(\C)$.  Hence $\phi$ is a surjective linear
endomorphism of the nine-dimensional vector space $M_3(\C)$ and is therefore
invertible.  This proves the theorem.
\end{proof}

\section{Verification, limitations, and novelty check}

\begin{itemize}[leftmargin=1.5em]
  \item \textbf{Internal consistency.} The range calculation uses an explicit
  preimage under $\Omega_\nu$, and the limit calculation follows directly
  from the source formula for $\pi_t^\#$. No computational or asymptotic
  approximation is used.
  \item \textbf{Classification check.} The three excluded nonfactor ranges
  are respectively $\C^3$, $\C\oplus M_2(\C)$, and $\C^2$; the two surviving
  factor ranges are $\C$ and $M_3(\C)$.
  \item \textbf{Scope.} The proof settles the full $n=3$ subcase only. It
  provides no classification for $n\geq4$.
  \item \textbf{Bounded search.} The four lightweight run indexes, the exact
  conjecture wording, the phrases ``q-pure $M_3$'' and ``unital q-pure maps,''
  and the arXiv records/citation trail for 0905.2708, 1005.4404, 1106.2304,
  and 1807.09824 were checked. The explicit $n=2$ result and the later q-weight
  classification were found; no paper explicitly stating the $n=3$ matrix-map
  corollary was located. Novelty confidence is therefore moderate, and the
  result is labeled literature-implied rather than wholly new.
  \item \textbf{Human-review focus.} Verify that Lemma 4.3 of the source paper
  covers every q-weight subordinate under the standard CP-flow/q-weight
  correspondence. Once that order transfer is accepted, the remaining steps
  are formal consequences of the cited theorems.
\end{itemize}

\begin{thebibliography}{9}
\bibitem{Jankowski2009}
C.~Jankowski,
\emph{On type II$_0$ $E_0$-semigroups induced by boundary weight doubles},
arXiv:0905.2708v2. Source conjecture on PDF p.~37; boundary-weight construction,
Corollary 3.3 and Lemma 4.3.

\bibitem{Jankowski2010}
C.~Jankowski,
\emph{Unital q-positive maps on $M_2(\C)$ and a related $E_0$-semigroup result},
arXiv:1005.4404v1. See Theorem 3.9 and Corollary 3.10 on PDF p.~20,
and Proposition 4.5 on PDF p.~23.

\bibitem{JMP2018}
C.~Jankowski, D.~Markiewicz, and R.~T.~Powers,
\emph{Classification of q-pure q-weight maps over finite dimensional Hilbert spaces},
arXiv:1807.09824v1. See Definition 1.2 and Theorem 5.7 on PDF p.~44.
\end{thebibliography}

\end{document}
